\documentclass[11pt]{amsart}
\usepackage[margin=1.1in]{geometry}
\usepackage{amsmath,amssymb,amsthm}
\usepackage{booktabs}
\usepackage{enumitem}
\usepackage[hidelinks]{hyperref}

\newtheorem{theorem}{Theorem}[section]
\newtheorem{lemma}[theorem]{Lemma}
\newtheorem{proposition}[theorem]{Proposition}
\newtheorem{corollary}[theorem]{Corollary}
\theoremstyle{definition}
\newtheorem{definition}[theorem]{Definition}
\newtheorem{example}[theorem]{Example}
\theoremstyle{remark}
\newtheorem{remark}[theorem]{Remark}

\newcommand{\F}{\mathbb{F}}
\newcommand{\Aut}{\operatorname{Aut}}
\newcommand{\gon}{\operatorname{gon}}
\newcommand{\Pl}{\mathcal{P}}
\newcommand{\minpoly}{\operatorname{mp}}
\newcommand{\Enc}{\operatorname{Enc}}
\newcommand{\code}[1]{\texttt{#1}}

\title[Orbits on places of small degree]{Orbits of $\Aut_{\F_5}(X_1(42))$ on places of small degree:\\ mathematics and computational details}
\author{Sourajyoti Maiti}
\address{Department of Mathematics, University of Connecticut}
\email{sourajyoti.maiti@uconn.edu}
\date{\today}

\begin{document}
\sloppy\emergencystretch=3em

\begin{abstract}
This note documents the computation of the orbits of the automorphism group $\Aut_{\F_5}(X)=\langle\sigma_4\rangle$ of $X=X_1(42)/\F_5$ on the places of degree $1,\dots,5$, and explains how the result is stored and reused. Every step is exact. The image of a place is found from the pull-back $\sigma^*$ alone, with no need to compute the inverse automorphism or any point map. The computation is split into independent parallel jobs, whose results are merged and checked. The note is a companion to the scripts \code{X1\_42\_orbits.m}, \code{X1\_42\_orbits\_merge.py}, \code{run\_orbits.sh}, \code{X1\_42\_orbits\_load.m} and \code{X1\_42\_orbits\_check.m}; it is not part of the manuscript.
\end{abstract}

\maketitle
\tableofcontents

%======================================================================
\section{Setup}
%======================================================================

\subsection{The curve}
Let $X=X_1(42)$ over $\F_5$; it has genus $25$. We use the plane model from Sutherland's tables (file \code{FFFc42.txt}). Its function field is
\[
K=\F_5(X)=\F_5(t)[y]/(F(t,y)),\qquad [K:\F_5(t)]=16,
\]
realised in Magma as \code{AFF := AlgorithmicFunctionField(FunctionField(XF2))}, with $t=$ \code{BaseField(AFF).1} and $y=$ \code{AFF.1}. All places are places of $K/\F_5$. For $v\ge1$ let $\Pl_v$ be the list \code{Places(AFF,v)} of places of degree $v$, in Magma's order, and write $P^{(v)}_i$ for its $i$-th entry. The sizes are
\[
\begin{array}{c|ccccccccc}
v & 1&2&3&4&5&6&7&8&9\\\hline
\#\Pl_v & 12&6&8&183&648&2288&11160&48708&214384
\end{array}
\]
The present computation covers $v\le5$.

\subsection{The group}
By the computation recorded in \code{X1\_42\_automorphisms.tex}, $G=\Aut_{\F_5}(X)$ is cyclic of order $12$, generated by $\sigma=\sigma_4$. The group is stored in \code{X1\_42\_aut\_F5.txt} as the pairs $(\sigma^*t,\sigma^*y)\in K^2$ of pull-backs. The scripts never recompute the group: they read the file with \code{LoadAutDataOpt(name,false)}, which takes about $2$\,s and skips the closure test. They then check that $\sigma_4$ has order exactly $12$, by composing $\sigma_4^*$ with itself until the identity $(t,y)$ appears; this takes about $87$\,s. On the rational places $\sigma_4$ acts as the permutation
\begin{equation}\label{eq:A4}
A_4=[9,5,2,1,12,8,11,7,3,4,10,6],
\end{equation}
which is the $12$-cycle $(1\,9\,3\,2\,5\,12\,6\,8\,7\,11\,10\,4)$. This was found independently by \code{X1\_42\_check\_orbs.m}.

%======================================================================
\section{Automorphisms acting on places}
%======================================================================

\subsection{From the pull-back to places}
Let $\phi\in G$ and write $\phi^*\colon K\to K$ for the corresponding $\F_5$-automorphism of the function field. For a place $P$ with valuation ring $\mathcal{O}_P$, put
\[
\mathcal{O}_{\phi(P)}=\{f\in K:\ \phi^*f\in\mathcal{O}_P\}=(\phi^*)^{-1}(\mathcal{O}_P).
\]
Since $\phi^*$ is a field automorphism fixing $\F_5$, this is again a valuation ring of $K/\F_5$. It defines a place $\phi(P)$.

\begin{lemma}\label{lem:basic}
For every place $P$ and every $f\in K$:
\begin{enumerate}[label=\textup{(\alph*)}]
\item $v_{\phi(P)}(f)=v_P(\phi^*f)$;
\item $\phi^*$ induces an $\F_5$-isomorphism $\bar\phi\colon\kappa(\phi(P))\to\kappa(P)$ of residue fields, with $\bar\phi(f(\phi(P)))=(\phi^*f)(P)$ whenever $f\in\mathcal{O}_{\phi(P)}$. In particular $\deg\phi(P)=\deg P$;
\item $P\mapsto\phi(P)$ is a bijection on the places of each degree $v$, and $\phi\mapsto(P\mapsto\phi(P))$ is an action of $G$, possibly through the opposite group, which does not change the orbits.
\end{enumerate}
\end{lemma}

\begin{proof}
(a) holds because $v_P\circ\phi^*$ is a normalised discrete valuation of $K$ with valuation ring $\mathcal{O}_{\phi(P)}$. For (b), $\phi^*$ maps $\mathcal{O}_{\phi(P)}$ onto $\mathcal{O}_P$ and, by (a), maps the maximal ideal onto the maximal ideal. It fixes $\F_5$, so the induced map of residue fields is an $\F_5$-isomorphism, and the degree is preserved. For (c), $(\phi^{-1})^*$ provides the inverse map, and $(\phi\psi)^*=\psi^*\phi^*$ shows that composition is respected up to order.
\end{proof}

Because $G=\langle\sigma\rangle$ is cyclic, the $G$-orbit of a place $P$ is $\{\sigma^k(P):0\le k<12\}$. Let $\pi_v$ be the permutation of $\{1,\dots,\#\Pl_v\}$ with $\sigma(P^{(v)}_i)=P^{(v)}_{\pi_v(i)}$. Then:

\begin{corollary}\label{cor:cycles}
The $G$-orbits on $\Pl_v$ are exactly the cycles of $\pi_v$. Every cycle length divides $12$, and $\pi_v^{12}=\mathrm{id}$.
\end{corollary}

So it suffices to compute the single permutation $\pi_v$ for each degree $v$. The difficulty is that $\sigma$ is known only through $\sigma^*$. Computing $\sigma(P)$ directly as a place would require the ideal $(\sigma^*)^{-1}(\mathcal{O}_P)$. Instead, we \emph{identify} $\sigma(P)$ inside the already computed list $\Pl_v$, using only valuations and residues of explicit functions at $P$. By Lemma~\ref{lem:basic}, these are values at $\sigma(P)$ in disguise.

%======================================================================
\section{Identifying the image of a place}
%======================================================================

Fix a degree $v$ and a place $P\in\Pl_v$. The image $\sigma(P)$ is found in at most three exact steps. Each step shrinks a set of candidates that is \emph{guaranteed} to contain $\sigma(P)$.

\subsection{Encoding polynomials}
For a polynomial $g=\sum_{i=0}^{n}c_ix^i\in\F_5[x]$ with $c_n\neq0$, put
\[
\Enc(g)=\sum_{i=0}^{n}c_i5^i+5^{n+1}\in\mathbb{Z}_{>0},\qquad c_i\in\{0,\dots,4\}.
\]
$\Enc$ is injective: the leading power $5^{n+1}$ recovers $n$, and the base-$5$ digits below it recover the coefficients. For $\alpha$ in a finite extension of $\F_5$, let $\minpoly(\alpha)\in\F_5[x]$ denote its minimal polynomial over $\F_5$.

\subsection{Step 1: the bucket}
The place of $\F_5(t)$ under a place $Q$ of $K$ is determined by $\minpoly(t(Q))$ if $t$ is regular at $Q$, and is the infinite place if $t$ has a pole at $Q$. Define the \emph{bucket} of $Q$ as
\[
b(Q)=\begin{cases}\Enc(\minpoly(t(Q))) & v_Q(t)\ge0,\\ 0 & v_Q(t)<0.\end{cases}
\]
We compute $b(Q)$ for all $Q\in\Pl_v$ once, and group the indices by bucket.

\begin{lemma}\label{lem:bucket}
Put $b^\sigma(P)=\Enc(\minpoly((\sigma^*t)(P)))$ if $v_P(\sigma^*t)\ge0$, and $b^\sigma(P)=0$ otherwise. Then $b(\sigma(P))=b^\sigma(P)$.
\end{lemma}

\begin{proof}
By Lemma~\ref{lem:basic}(a), $v_{\sigma(P)}(t)=v_P(\sigma^*t)$, so $t$ has a pole at $\sigma(P)$ iff $\sigma^*t$ has a pole at $P$. Otherwise $\bar\sigma(t(\sigma(P)))=(\sigma^*t)(P)$ by Lemma~\ref{lem:basic}(b). An $\F_5$-isomorphism of fields preserves minimal polynomials over $\F_5$.
\end{proof}

So the first candidate set is $C_1=\{j:\ b(P^{(v)}_j)=b^\sigma(P)\}\subseteq\Pl_v$. It contains $\pi_v(i)$ for $P=P^{(v)}_i$, and has at most $16$ elements, since $\sum e_if_i=[K:\F_5(t)]=16$ over any place of $\F_5(t)$. If $\#C_1=1$, the image is found.

\subsection{Step 2: the key}
For a place $Q$ at which $t$ and $y$ are both regular, put $a=t(Q)$ and $c=y(Q)$, and define its \emph{key}
\[
\kappa(Q)=\sum_{m=0}^{3}\Enc(g_m)\,(5^{12})^{3-m},\qquad (g_0,g_1,g_2,g_3)=\bigl(\minpoly(a),\minpoly(c),\minpoly(a+c),\minpoly(a+2c)\bigr).
\]
If $t$ or $y$ has a pole at $Q$, we set $\kappa(Q)=0$, meaning ``not available''. All $g_m$ have degree at most $v\le9$, so $\Enc(g_m)<5^{11}<5^{12}$, and $\kappa$ is an injective function of the quadruple $(g_0,\dots,g_3)$.

\begin{lemma}\label{lem:key}
If $\sigma^*t$ and $\sigma^*y$ are regular at $P$, then $\kappa(\sigma(P))$ equals the key computed from the pair $\bigl((\sigma^*t)(P),(\sigma^*y)(P)\bigr)$ in place of $(a,c)$.
\end{lemma}

\begin{proof}
Use the same isomorphism $\bar\sigma$ for both coordinates. It sends $(t(\sigma(P)),y(\sigma(P)))$ to $((\sigma^*t)(P),(\sigma^*y)(P))$, and hence sends $a+c$ and $a+2c$ to the corresponding combinations. Minimal polynomials over $\F_5$ are preserved.
\end{proof}

The second candidate set is $C_2=\{j\in C_1:\ \kappa(P^{(v)}_j)=\kappa^\sigma(P)\}$. It still contains the true image. If $\#C_2=1$, the image is found. If $C_2$ is empty (which can only happen when a key is unavailable), or if $\kappa^\sigma(P)$ cannot be computed, we continue with $C=C_1$; otherwise we continue with $C=C_2$.

\begin{remark}
Two different places can have the same key: for example, two places over the same place of $\F_5(t)$ whose residues $(a,c)$ are conjugate over $\F_5$ only coordinatewise, but not as a pair. The four minimal polynomials make such collisions rare, and the step is cheap. Step 3 handles whatever remains.
\end{remark}

\subsection{Step 3: the uniformizer pattern}
Let $C=\{j_1,\dots,j_r\}$ with $r\ge2$, and write $Q_s=P^{(v)}_{j_s}$. For every candidate $Q_s$, choose a function $u_s\in K$ with $v_{Q_s}(u_s)>0$. In practice $u_s=\code{LocalUniformizer}(Q_s)$, and the script checks $v_{Q_s}(u_s)>0$. Nothing is required of $u_s$ at the other candidates. Record:
\begin{itemize}
\item the zero pattern $Z_s=\{q\in\{1,\dots,r\}:\ v_{Q_q}(u_s)>0\}$, which always contains $s$;
\item the test bit $\beta_s=\bigl[\,v_P(\sigma^*u_s)>0\,\bigr]\in\{0,1\}$.
\end{itemize}
A candidate $q$ \emph{survives} if $q\in Z_s\iff\beta_s=1$ for every $s=1,\dots,r$.

\begin{proposition}\label{prop:unif}
Let $q^\ast$ be the index with $Q_{q^\ast}=\sigma(P)$.
\begin{enumerate}[label=\textup{(\alph*)}]
\item $q^\ast$ always survives.
\item If exactly one candidate survives, it is $q^\ast$.
\item A candidate $q\neq q^\ast$ survives only if $u_q$ vanishes at $\sigma(P)$.
\end{enumerate}
\end{proposition}

\begin{proof}
By Lemma~\ref{lem:basic}(a), $\beta_s=[v_{\sigma(P)}(u_s)>0]=[q^\ast\in Z_s]$, which proves (a). (b) follows from (a). For (c): $q\in Z_q$, so $q$ can survive only if $\beta_q=1$, i.e. $v_{\sigma(P)}(u_q)>0$.
\end{proof}

So Step 3 is never wrong; at worst it is inconclusive. In that case the image is recorded as $-1$ (\emph{ambiguous}), and the whole degree is rejected at the merge stage (Section~\ref{sec:merge}). Ambiguity is unlikely: by (c), it requires the chosen uniformizer of another candidate to vanish at the true image. Should it ever occur, it can be removed by replacing $u_q$ by another function vanishing at $Q_q$ but not at the competing candidates (e.g.\ $u_q+\lambda u_q^2$ or a weak-approximation element), or by the constraint that $\pi_v$ is a bijection.

\begin{example}
Suppose $C=\{1,2,3\}$, with zero patterns $Z_1=\{1\}$, $Z_2=\{2,3\}$ and $Z_3=\{3\}$, and observed bits $(\beta_1,\beta_2,\beta_3)=(0,1,1)$. Candidate $1$ fails at $s=1$, since $1\in Z_1$ but $\beta_1=0$. Candidate $2$ fails at $s=3$, since $2\notin Z_3$ but $\beta_3=1$. Candidate $3$ satisfies all three conditions. Hence $\sigma(P)=Q_3$. With bits $(0,1,0)$ instead, only candidate $2$ survives.
\end{example}

\subsection{Cost}
The functions $\sigma^*t$, $\sigma^*y$ and $\sigma^*u_s=u_s(\sigma^*t,\sigma^*y)$ (\code{ApplyAut}) are exact elements of $K$. Steps 1 and 2 need two residue evaluations and four minimal polynomials per place. Step 3 needs at most $16$ local uniformizers, $16^2$ valuations at the candidates, and $16$ valuations of pulled-back functions at $P$. The pulled-back functions are the expensive part, which is why Step 3 is used only when Steps 1--2 leave several candidates. The script reports, for each degree, how many images were decided by each step.

%======================================================================
\section{Checks}\label{sec:checks}
%======================================================================

The following are verified before a degree is saved.
\begin{enumerate}[label=(C\arabic*)]
\item \emph{Group.} The loaded $\sigma_4$ has order exactly $12$.
\item \emph{Consistent numbering.} All parallel jobs for a degree use the same list $\Pl_v$. Each job writes $\#\Pl_v$ and the checksum
\[
\mathrm{CHK}=\Bigl(\cdots\bigl((0\cdot B+h_1)\,B+h_2\bigr)\cdots\Bigr)\bmod M,\qquad h_i=\bigl(b_i\bmod p\bigr)+7\bigl(\kappa_i\bmod p\bigr),
\]
with $B=1000003$, $p=10^9+7$, $M=2^{61}-1$, and $b_i,\kappa_i$ the bucket and key of $P^{(v)}_i$. The merge requires all parts to agree.
\item \emph{Completeness.} Every index has exactly one image, computed by the job responsible for it, and no image is ambiguous.
\item \emph{Permutation.} The images form a permutation of $\{1,\dots,\#\Pl_v\}$. By Lemma~\ref{lem:basic}(c), this must hold, so it is a non-trivial check of the whole identification.
\item \emph{Group law.} All cycle lengths divide $12$ (Corollary~\ref{cor:cycles}).
\item \emph{Known case.} For $v=1$, the permutation equals $A_4$ in \eqref{eq:A4} or its inverse. The inverse is accepted because the convention $\phi\mapsto\phi(P)$ versus $\phi\mapsto\phi^{-1}(P)$ only replaces $\pi_v$ by $\pi_v^{-1}$, with the same orbits.
\end{enumerate}
If any check fails for a degree, no file is written for that degree, and an old file of that degree is deleted, so a stale result can never be reused.

%======================================================================
\section{Parallel computation}\label{sec:merge}
%======================================================================

\subsection{Splitting}
The images of different places are independent, so the work splits trivially. With $m$ parts for degree $v$, job $k\in\{0,\dots,m-1\}$ computes $\pi_v(i)$ exactly for
\[
i\equiv k+1\pmod m.
\]
This round-robin split balances the load better than contiguous blocks, because Magma's order of $\Pl_v$ groups places by the place of $\F_5(t)$ below them, and bucket sizes vary.

Each job still computes the full list $\Pl_v$ and the buckets and keys of \emph{all} places, since any of them may be a candidate. It also loads the group and checks the order. These shared costs are small compared with the image computations for $v=4,5$.

\subsection{Launch and merge}
\code{run\_orbits.sh P1 P2 P3 P4 P5} starts $P_v$ jobs for degree $v$; the default is $1,1,1,4,12$, i.e.\ $19$ jobs, one per core:
\[
\code{magma -b degs:=}v\ \code{part:=}k\ \code{nparts:=}m\ \code{X1\_42\_orbits.m}
\]
Job $(v,k,m)$ writes \code{X1\_42\_orbits\_d<v>\_part<k>\_of\_<m>.txt}, which contains:
\begin{itemize}
\item a header with \code{DEG}, \code{GEN}, \code{PART}, \code{NPARTS}, \code{COUNT} and \code{CHECKSUM};
\item for part $0$ only, the place data lines \code{P i | c0 c1 ... | key};
\item its lines \code{I i j}, meaning $\pi_v(i)=j$;
\item a closing \code{END} line.
\end{itemize}
After all jobs finish, \code{X1\_42\_orbits\_merge.py} performs checks (C2)--(C6), splits $\pi_v$ into cycles, and writes the final files.

%======================================================================
\section{The saved data}
%======================================================================

For each degree $v$, the file \code{X1\_42\_orbits\_d<v>.txt} contains:
\begin{center}\small
\begin{tabular}{@{}lp{0.5\textwidth}@{}}
\toprule
line & meaning\\\midrule
\code{X1\_42\_ORBITS v1 DEG v GEN 4} & header, continued by\\
\quad\code{COUNT n ORBITS k CHECKSUM c} & \\
\code{P i | c0 c1 ... cn | key} & place $i$: coefficients of $\minpoly(t(P_i))$, constant term first (empty for $t=\infty$), and $\kappa(P_i)$\\
\code{I i j} & $\sigma_4(P_i)=P_j$\\
\code{O k s : i1 i2 ... is} & orbit $k$ of size $s$, as the cycle $i_1\mapsto i_2\mapsto\cdots$\\
\code{R k i} & representative of orbit $k$ (its smallest index)\\
\code{END n} & end marker\\\bottomrule
\end{tabular}
\end{center}
\code{X1\_42\_orbits\_summary.txt} lists, for each degree, the number of orbits and the distribution of orbit sizes.

\subsection{Reuse}
\code{X1\_42\_orbits\_load.m} provides three functions.
\begin{itemize}
\item \code{LoadOrbitFile(v)} reads the file into a record with fields \code{coeffs}, \code{keys}, \code{img}, \code{orbits}, \code{reps} and \code{orbitof}. It needs no curve computation.
\item \code{PlacesFromOrbitFile(D)} recomputes $\Pl_v$ and verifies, place by place, the stored polynomial and key and the checksum. Afterwards, index $i$ in the file is guaranteed to refer to \code{PL[i]} in the current session.
\item \code{FindPlace(coeffs, key, v)} recovers a single place without computing $\Pl_v$. It takes the zeros of $p(t)$, where $p$ is the stored polynomial (or the poles of $t$ if the data is empty), keeps those of degree $v$, and keeps the one with matching key. It succeeds when the stored data singles out the place; the merge reports how many places do not satisfy this.
\end{itemize}
\code{X1\_42\_orbits\_check.m} runs both reloading routes on every saved degree and prints PASS/FAIL lines.

%======================================================================
\section{Use in the divisor search}
%======================================================================

Let $S$ be the sum of the $12$ rational places (the cusps). Since $G$ permutes the rational places, $\phi(S)=S$ for all $\phi\in G$. The manuscript (Corollary~3.3) reduces the statement $\gon_{\F_5}(X)\ge d+1$ to the following: for every effective divisor $E$ of degree $d-2$ in a certain family $\mathcal{E}$, every nonconstant $h\in L(S+E)$ has degree $>d$. For $\phi\in G$, the map $h\mapsto h\circ\phi^{-1}$ is a degree-preserving bijection $L(S+E)\to L(S+\phi(E))$ (Lemma~4.1 of the manuscript). Hence it suffices to test one $E$ from every $G$-orbit, provided $\mathcal{E}$ is $G$-stable.

Represent $E$ by the multisets of indices of its places in each $\Pl_v$. Let $\mathsf{K}$ be the set of degrees for which $\pi_v$ is known, and let $E_{\mathsf{K}}$ be the part of $E$ supported in those degrees.

\begin{proposition}\label{prop:canon}
Order the representations of $E_{\mathsf{K}}$ lexicographically, and \emph{keep} $E$ if and only if $E_{\mathsf{K}}\le\sigma^k(E_{\mathsf{K}})$ for $k=1,\dots,11$. If $\mathcal{E}$ is $G$-stable, every $G$-orbit in $\mathcal{E}$ contains a kept divisor. If every degree occurring in divisors of $\mathcal{E}$ lies in $\mathsf{K}$, every orbit contains exactly one kept divisor.
\end{proposition}

\begin{proof}
Given $E$, choose $k$ minimising $\sigma^k(E)_{\mathsf{K}}=\sigma^k(E_{\mathsf{K}})$, and put $E'=\sigma^k(E)\in\mathcal{E}$. For every $j$, $\sigma^j(E'_{\mathsf{K}})=\sigma^{j+k}(E_{\mathsf{K}})\ge E'_{\mathsf{K}}$, so $E'$ is kept. Note that the images of places of degree outside $\mathsf{K}$ are never needed. If all degrees lie in $\mathsf{K}$, then $E_{\mathsf{K}}=E$, and the lexicographic minimum of an orbit is unique.
\end{proof}

Degrees without a known permutation are thus simply not reduced, which is always safe. For $d=11$ a test with synthetic permutations of the real place counts reduced $2{,}994{,}534$ divisors to $531{,}211$, a factor of about $5.6$; the factor is below the group order $12$ because divisors with nontrivial stabiliser have smaller orbits, and the parts of degree outside $\mathsf{K}$ are not reduced. The reduction is implemented in \code{X1\_42\_symreduce.py} (stage B). The search scripts (stage C) check, before searching, that the place numbering in the session agrees with the stored keys and checksum.

%======================================================================
\section{Status and extension}
%======================================================================

\begin{itemize}
\item \emph{Tested commands.} \code{LocalUniformizer} returned a function with valuation $1$ at all $20$ test places (test T8b of \code{X1\_42\_symtest.m}). \code{UniformizingParameter} and \code{WeakApproximation([Q],[1])} are not available with these argument types in Magma V2.29-10, and are not used, except that \code{WeakApproximation} is tried first inside a \code{try}-block in the search scripts. Loading the group takes about $2$\,s, and the order check about $87$\,s.
\item \emph{To be confirmed on the server.} The output of \code{run\_orbits.sh}, the merge summary, the PASS lines of \code{X1\_42\_orbits\_check.m}, and test T11, which compares degrees $1$--$3$ against $A_4$.
\item \emph{Higher degrees.} The method applies unchanged to $v=6,\dots,9$; only more parts are needed. Computing $\Pl_9$ (with $214{,}384$ places) is itself a substantial computation, repeated in every job. For these degrees, the place list and the data of Steps 1--2 should therefore be computed once and stored before the image jobs start.
\end{itemize}

\end{document}
